% 351_PDG_Modellabhaengigkeit_De_ch.tex
% Johann Pascher, 5. September 2026

\providecommand{\BM}{\textbf{[B]}}
\providecommand{\KM}{\textbf{[K]}}
\providecommand{\SM}{\textbf{[S]}}
\providecommand{\QM}{\textbf{[Q]}}

\phantomsection
\addcontentsline{toc}{chapter}{Modellabhängigkeit der PDG-Vergleichswerte}

\section*{Statuszeichen}
\begin{center}
\begin{tabular}{@{}lp{10.5cm}@{}}
\textbf{[K]} & \emph{Konfirmiert} --- aus $\xi$ hergeleitet, numerisch geprüft. \\[4pt]
\textbf{[B]} & \emph{Belegt} --- algebraisch bewiesen. \\[4pt]
\textbf{[S]} & \emph{Spekulativ} --- offen. \\[4pt]
\textbf{[Q]} & \emph{Quelle} --- korpusexterne Primärquelle, hier verifiziert. \\
\end{tabular}
\end{center}

\section*{Abstract}

FFGFT-Vorhersagen werden gegen PDG- und CODATA-Werte verglichen. Dieses Dokument untersucht, wie diese Vergleichswerte tatsächlich zustande kommen: welche Rohobservable zugrunde liegt, welche Theorie (QED, Standardmodell) in die Extraktion eingeht, und wo irrationale Zahlen bereits in den Extraktionsformeln auftreten.

Hauptbefunde: (1)~$m_e$ und $m_\mu/m_e$ sind die modellärmsten Anker, aber $m_\mu/m_e$ ist QED-abhängig und stammt aus einer 27~Jahre alten Myonium-HFS-Messung, deren Theorieunsicherheit CODATA laut Eides (2026) um Faktor 5--10 unterschätzt. (2)~$\alpha$ hat eine ungelöste $5\sigma$-Spannung zwischen den beiden besten Rückstoßmessungen. (3)~$v=246$~GeV wird nie gemessen; es ist die SM-Definition $v\equiv(\sqrt2\,G_F)^{-1/2}$. (4)~$G_F$ trägt eine Theoriekorrektur $\Delta q=-0{,}44\,\%$ mit hadronischer $4{,}7\sigma$-Spannung. (5)~$m_\tau$ hat zwei Messmethoden, die um 0{,}18~MeV streuen; FFGFT liegt $0{,}4\sigma$ vom PDG-Mittel. (6)~Die Größenordnung aller Theorieunsicherheiten ($10^{-7}$ bis $10^{-10}$) erklärt das $\sim1\,\%$-Residuum bei $(m_\mu/m_e)^2$ nicht --- es bleibt physikalisch. (7)~Fehler akkumulieren entlang der Ableitungskette; jede Vorhersage braucht ihren vollständigen Herkunftspfad.

Alle 21 Primärquellen wurden am 5.~September 2026 verifiziert.

\section{Fragestellung}

Der Corpus vergleicht FFGFT-Vorhersagen mit PDG-Werten und findet Residuen im Sub-Prozent- bis Prozent-Bereich. Sind die PDG-Werte theorie-neutral, oder hängen sie selbst am Standardmodell --- sodass ein Teil der Residuen nicht FFGFT-seitig, sondern extraktionsseitig entsteht?

Die FFGFT-Grundaussage lautet: Nur Verhältnisse ohne irrationale Zahlen sind exakt; SI-Absolutwerte sind strukturell approximativ (Dok.~241, 085). Wenn die PDG-Werte selbst irrationale Faktoren in der Extraktion enthalten, gilt dieselbe Einschränkung für die Vergleichsseite.

\section{Wie empirische Messwerte entstehen}

Kein PDG- oder CODATA-Wert ist ein Rohmesswert. Jeder durchläuft dieselbe Kette:
\begin{center}
Rohobservable $\to$ Theorieformel $\to$ extrahierter Parameter $\to$ CODATA-Ausgleich $\to$ PDG-Wert
\end{center}

\subsection{Der CODATA-Ausgleich}

CODATA veröffentlicht alle vier Jahre einen Satz empfohlener Werte. Diese entstehen aus einem gemeinsamen Ausgleich nach der Methode der kleinsten Quadrate über etwa 300 Eingangsdaten, die über Theoriebeziehungen (QED, Rydberg-Formel, Kinematik) verknüpft sind. Die Werte sind daher \emph{korreliert}: Ändert sich ein Eingangsdatum, verschieben sich alle davon abhängigen Werte. Wo Theorieunsicherheiten in den Ausgleich eingehen, werden sie zum Teil als Ausgleichsparameter behandelt ($\delta_\text{th}$ in den Observationsgleichungen), nicht als Fehler --- das ist der Eides-Befund bei der Myonium-HFS \QM.

\subsection{Elektronmasse $m_e$}

\textbf{Rohobservable:} Zyklotronfrequenz eines wasserstoffartigen Ions ($^{12}$C$^{5+}$, $^{28}$Si$^{13+}$, $^{20}$Ne$^{9+}$) in einer Penning-Falle und Larmor-Präzessionsfrequenz des gebundenen Elektrons. Das Verhältnis ist eine reine Zahl.

\textbf{Theorieformel:} Der gebundene $g$-Faktor aus QED-Bound-State-Theorie. Aus Frequenzverhältnis und $g_\text{bound}$ folgt $m_e/m_\text{Ion}$, mit Ionenmasse in u folgt $m_e$ in u; Umrechnung in MeV über u\,$\to$\,kg\,$\to$\,MeV.

\textbf{Theorieabhängigkeit:} QED für $g_\text{bound}$, klein und gut kontrolliert. Aktuell: Heiße et al. (2023) auf 0{,}1~ppb \QM.

\subsection{Myon-Elektron-Massenverhältnis $m_\mu/m_e$}

\textbf{Rohobservable:} Hyperfeinaufspaltung des Myonium-Grundzustands --- Mikrowellenresonanz bei 4{,}463~GHz --- und Zeeman-Aufspaltung im Magnetfeld.

\textbf{Theorieformel:} $\Delta\nu_\text{HFS}=\nu_F\,[1+\text{Korrekturen}]$ mit der Fermi-Frequenz
\begin{equation}
\nu_F=\tfrac{16}{3}\,\alpha^2\,c\,R_\infty\,\frac{m_e}{m_\mu}\Big(\frac{m_r}{m_e}\Big)^3 .
\end{equation}
Korrekturen: QED bis $O(\alpha^3)$, Rückstoß in $m_e/m_\mu$, schwacher $Z$-Austausch, hadronische Vakuumpolarisation.

\textbf{Theorieabhängigkeit:} hoch. Einzige Präzisionsmessung: LAMPF 1999 (Liu et al.). Eides (2026) zeigt, dass CODATA den Messwert 4\,463\,302\,776(51)~Hz als theoretischen Wert führt; die tatsächliche Theorieunsicherheit beträgt 271--515~Hz \QM.

\textbf{Für FFGFT:} $\nu_F\propto\alpha^2$. Der Anker $m_\mu/m_e$ enthält $\alpha$ implizit. Anker und $\alpha$ sind nicht unabhängig.

\subsection{Feinstrukturkonstante $\alpha$}

Zwei Wege, die sich um $5\sigma$ widersprechen \QM:
\begin{itemize}\itemsep2pt
\item \textbf{Atomrückstoß:} $\alpha^2=2R_\infty\,(h/m_\text{Atom})\,(m_\text{Atom}/m_e)/c$. Rb (Morel et al.\ 2020): $\alpha^{-1}=137{,}035\,999\,206(11)$. Cs (Parker et al.\ 2018): $137{,}035\,999\,046(27)$. Differenz $1{,}6\times10^{-7}$, $>5\sigma$.
\item \textbf{Elektron $g-2$:} $a_e$ auf $10^{-13}$ (Fan et al.\ 2023) plus 5-Loop-QED (12\,672 Diagramme, Aoyama et al.). Ergibt $\alpha^{-1}=137{,}035\,999\,166(15)$ --- zwischen Rb und Cs.
\end{itemize}
Der CODATA-Wert 137{,}035\,999\,177(21) ist ein Mittel, das keine der Messungen abdeckt. FFGFT: $\alpha^{-1}=3700/27=137{,}037\,037\ldots$ (Dok.~338, R101), Abweichung 7{,}6~ppm --- 100$\times$ größer als die Rb/Cs-Spannung.

\subsection{Tau-Masse $m_\tau$ --- zwei Methoden}

\begin{center}
\begin{tabular}{@{}llll@{}}
\toprule
Methode & Experiment & Wert (MeV) & Theorieannahme \\
\midrule
Schwellenscan & BESIII 2014 & $1776{,}91\pm0{,}12\,{}^{+0{,}10}_{-0{,}13}$ & QED-Wirkungsquerschnitt \\
Pseudomasse & Belle II 2023 & $1777{,}09\pm0{,}08\pm0{,}11$ & $m_{\nu_\tau}=0$ \\
\textbf{PDG-Mittel 2024} & --- & $\mathbf{1776{,}93\pm0{,}09}$ & --- \\
\bottomrule
\end{tabular}
\end{center}

\textbf{Schwellenscan:} $\sigma(e^+e^-\to\tau^+\tau^-)$ gegen Strahlenergie nahe $2m_\tau\approx3{,}55$~GeV; $m_\tau$ als Fitparameter der QED-Schwellenkurve. \textbf{Pseudomasse:} Bei 10{,}6~GeV aus $\tau\to3\pi\nu$ die kinematische Kante der invarianten Masse; setzt $m_{\nu_\tau}=0$ voraus.

\textbf{Für FFGFT:} $m_\tau=1776{,}97$~MeV (pred) liegt 0{,}04~MeV $=0{,}4\sigma$ vom PDG-Mittel. Die Methoden streuen um 0{,}18~MeV. Der Falsifikationstest ist nicht entschieden; Belle II mit $>1$~ab$^{-1}$ sollte $\sigma<0{,}05$~MeV erreichen.

\subsection{Fermi-Konstante $G_F$ und Vakuumerwartungswert $v$}

\textbf{Rohobservable für $G_F$:} Zeitverteilung des Myonzerfalls (MuLan, $10^{12}$ Zerfälle), $\tau_\mu=2{,}196\,980\,3(22)$~µs \QM.

\textbf{Theorieformel} (Eberhart et al.\ 2026 \QM):
\begin{equation}
\tau_\mu^{-1}=\frac{G_F^2\,m_\mu^5}{192\pi^3}\,(1+\Delta q),\qquad \Delta q=(-4\,384\,678\pm34)\times10^{-9}\approx-0{,}44\,\%.
\end{equation}
$\Delta q$ enthält QED bis $O(\alpha^3)$, hadronische Vakuumpolarisation, Elektronmasseneffekte. Rein elektroschwache Korrekturen sind per Definition in $G_F$ absorbiert (Sirlin 1980). Die hadronische Komponente hat eine $4{,}7\sigma$-Spannung (CMD-3 vs.\ ältere Daten).

\textbf{Rohobservable für $v$:} keine. $v\equiv(\sqrt2\,G_F)^{-1/2}=246{,}22$~GeV folgt aus $M_W=gv/2$ und $G_F/\sqrt2=g^2/(8M_W^2)$. \textbf{$v$ existiert nur im Standardmodell.} Der Zahlenwert ist die Myon-Lebensdauer in SM-Kleidung.

\textbf{Für FFGFT:} Dok.~006/R104 gibt $v=248{,}3$~GeV (ohne $K_\text{frak}$) bzw.\ 245{,}6~GeV (mit $K_\text{frak}=74/75$). Der Vergleich mit 246{,}22~GeV ist ein Vergleich \emph{FFGFT-Struktur gegen SM-Definition}, nicht gegen ein Observable. Nur $\tau_\mu$ und $m_\mu$ sind gemessen.

\subsection{Gravitationskonstante $G$}

Torsionswaage (Li et al.\ 2018, 11{,}6~ppm) oder Atominterferometrie (Rosi et al.\ 2014, 150~ppm). Theorie: Newton. Einziger modellfreier Wert, aber der ungenaueste: CODATA $6{,}67430(15)\times10^{-11}$, 22~ppm \QM.

\subsection{Neutrino-Massendifferenzen}

Zählraten gegen $L/E$, Dreiflavour-Fit (NuFIT 6.0 \QM). $\Delta m^2_{31}$ hängt davon ab, ob SK-Atmosphärendaten einbezogen werden: Verhältnis 33{,}7 (mit) vs.\ 33{,}9 (ohne). FFGFT $360/11=32{,}73$; Messunsicherheit 1--3\,\% $\approx$ FFGFT-Rest. Nicht diskriminierend.

\subsection{Übersicht}

\begin{center}
\begin{tabular}{@{}llll@{}}
\toprule
Größe & Rohobservable & Theorie in Extraktion & Konvention \\
\midrule
$m_e$ & Frequenzverhältnis & QED ($g_\text{bound}$) & --- \\
$m_\mu/m_e$ & HFS-Frequenz & QED, schwach, hadr. & --- \\
$\alpha$ (Rb/Cs) & Interferometerphase & $R_\infty$ (QED) & --- \\
$\alpha$ ($g-2$) & Frequenzverhältnis & 5-Loop-QED & --- \\
$m_\tau$ & Zählraten vs.\ $E$ / Kante & QED-Schwelle / $m_{\nu_\tau}=0$ & Strahlkalibration \\
$G_F$ & Zerfallszeiten & Fermi + QED + hadr. & Sirlin-Absorption \\
$v$ & --- & --- & $v\equiv(\sqrt2G_F)^{-1/2}$ \\
$G$ & Winkel / Beschleunigung & Newton & --- \\
$\Delta m^2$ & Zählraten vs.\ $L/E$ & 3-Flavour & mit/ohne SK-atm \\
\bottomrule
\end{tabular}
\end{center}

Kein Wert ist theoriefrei. Für FFGFT bedeutet das: der Vergleich mit PDG-Werten ist ein Vergleich mit \glqq QED + Konvention\grqq, nicht mit rohen Beobachtungen.

\section{Irrationale Faktoren in den Extraktionsformeln}

Die FFGFT-Grundaussage --- Verhältnisse rationaler Zahlen sind exakt, SI-Absolutwerte tragen die Approximation der Projektion $S^1_m\to\mathbb R_t$ (Dok.~306, 333) --- gilt für die SM-Extraktionsformeln selbst:

\begin{center}
\begin{tabular}{@{}lll@{}}
\toprule
Größe & Formel & Irrationaler Faktor \\
\midrule
$G_F$ & $\tau_\mu^{-1}=G_F^2m_\mu^5/(192\pi^3)$ & $\pi^3$ \\
$v$ & $v=(\sqrt2\,G_F)^{-1/2}$ & $\sqrt2$ \\
$\alpha$ (Rückstoß) & $\alpha^2=2R_\infty h/(mc)$ & $R_\infty\propto\alpha^2$: implizit zirkulär \\
HFS & $\nu_F\propto\alpha^2$ & $\alpha^2$ \\
$\Delta q$ & Entwicklung in $\alpha/\pi$ & $\pi$ in jedem Term \\
\bottomrule
\end{tabular}
\end{center}

\glqq$G_F$ gemessen\grqq{} heißt: $\tau_\mu$ gemessen, durch $192\pi^3$ geteilt, Wurzel gezogen. \glqq$v$ gemessen\grqq{} heißt zusätzlich: durch $\sqrt2$ geteilt, nochmals Wurzel. Die irrationalen Faktoren stammen aus der SM-Formulierung (Phasenraumintegral, Higgs-Normierung), nicht aus der Messung. Der einzige konventionsfreie Vergleich wäre auf der Ebene von $\tau_\mu=2{,}196\,980\,3$~µs --- was FFGFT direkt vorhersagen müsste \SM.

\section{Das Residuum bei $(m_\mu/m_e)^2$: exakte Rechnung}

FFGFT bare (Dok.~338): $(m_\mu/m_e)^2=43200$ exakt, Integer, $=|\text{GF}(9)^*|^2\cdot5^2\cdot|\text{GF}(27)|$. Mit CODATA 2022 in exakter rationaler Arithmetik:
\begin{align}
(m_\mu/m_e)^2_\text{PDG} &= 42753{,}12 \\
\text{Rest} &= 446{,}88\quad(1{,}045\,\%) \\
\sigma_\text{rel}\big((m_\mu/m_e)^2\big) &= 4{,}35\times10^{-8}\quad\Rightarrow\quad \text{Rest}=240\,000\,\sigma .
\end{align}

Mögliche Fehlerquellen, durchgerechnet:
\begin{itemize}\itemsep2pt
\item CODATA-Unsicherheit: $\pm0{,}002$ --- Faktor $2\times10^5$ zu klein.
\item Eides-Korrektur der HFS-Theorieunsicherheit (515 statt 51~Hz): $\pm0{,}010$ --- Rest bleibt $45\,000\,\sigma$.
\item $\alpha$-Wahl im HFS-Anker (FFGFT $3700/27$ statt CODATA): verschiebt $(m_\mu/m_e)^2$ um $-1{,}3$ --- Faktor 500 zu klein.
\item Um 43200 zu erreichen, müsste $m_\mu$ um 551~keV höher sein (Messfehler 2{,}3~eV), oder $m_e$ um 2{,}65~keV niedriger (Messfehler 0{,}15~meV), oder die HFS um 23~MHz (Messfehler 51~Hz).
\end{itemize}

\textbf{Fazit \BM:} Kein Messfehler erklärt den Rest. Er ist physikalisch und gehört zur Theorie: fraktal-rekursive Korrektur (Dok.~295, 146, 338) und Projektionsapproximation (Dok.~306, 333). Im Quadrat, ohne Wurzel: $43200\cdot(1-77{,}6\,\xi)=42753$; $K_\text{frak}=1-100\xi$ überkorrigiert um $22{,}4\,\xi$. Welche Potenz von $K_\text{frak}$ auf die quadratische Identität wirkt, ist im Corpus nicht festgelegt \SM.

\section{Fehlerakkumulation entlang der Ableitungskette}

Jeder Schritt trägt einen kleinen Rest (Rekursion $\sim100\xi$, Projektion $\sim\xi$, Extraktionstheorie $10^{-7}$--$10^{-10}$). Schon bei $m_\mu$, dem ersten Schritt nach dem Anker $m_e$:

\begin{center}
\begin{tabular}{@{}ll@{}}
\toprule
Rechenweg $m_e\to m_\mu$ & Rest \\
\midrule
Galois bare: $m_\mu=\sqrt{43200}\,m_e$ & $+0{,}52\,\%$ \\
mit $\sqrt{K_\text{frak}}$: $m_\mu=\sqrt{43200\cdot74/75}\,m_e$ & $-0{,}15\,\%$ \\
über $\alpha$: $E_0=\sqrt{\alpha/\xi}$, $m_\mu=E_0^2/m_e$ & $+1{,}37\,\%$ \\
\bottomrule
\end{tabular}
\end{center}

Bei $m_\tau$ (zweiter Schritt), $v$ (dritter, über $\xi^{3/2}$) akkumuliert jede Ableitung weiter. Um zu entscheiden, wo ein Rest sitzt, müsste die Konvergenz der empirischen Werte auf die algebraischen gesichert sein --- das ist sie nicht: $m_\mu/m_e$ (1999), $\alpha$ ($5\sigma$-Spannung), $\Delta m^2_\text{atm}$ (2{,}4$\sigma$-Sprung 2024), $v$ (nicht gemessen).

\textbf{Buchführung:} Für jede berichtete Vorhersage sind anzugeben: (1)~der Anker, (2)~die Anzahl der Ableitungsschritte, (3)~angewendete Korrekturen an welcher Stelle, (4)~der Vergleichswert mit Extraktionsmethode. \glqq FFGFT sagt $m_\mu$ mit 0{,}5\,\% vorher\grqq{} ist unvollständig; vollständig: \glqq Von $m_e$ (PDG) aus, in einem Schritt über die Galois-Identität, ohne $K_\text{frak}$, ergibt sich $m_\mu$ mit $+0{,}52\,\%$ gegenüber dem Myonium-HFS-Wert von 1999.\grqq

\section{Konsequenzen für den Corpus}

\textbf{Kein Zahlenwert im Corpus ändert sich.} Was sich ändert:
\begin{enumerate}\itemsep2pt
\item \textbf{$v$ ist kein Anker und kein Observable.} Dok.~006, 319, R104: \glqq$v=246$~GeV (gemessen)\grqq{} $\to$ \glqq$v=246$~GeV (SM-Definition aus $G_F$)\grqq. Register R112.
\item \textbf{$m_\mu/m_e$ ist QED-abhängig.} \glqq modellunabhängig\grqq{} $\to$ \glqq modellärmster verfügbarer Wert\grqq; $\nu_F\propto\alpha^2$. Register R113.
\item \textbf{$\alpha$ hat keinen eindeutigen Messwert unter 1~ppb.} Vergleichswert als CODATA-Mittel mit Rb/Cs-Spannung zitieren.
\item \textbf{$m_\tau$:} zwei Methoden, PDG-Mittel 1776{,}93(9), FFGFT $0{,}4\sigma$ --- offen.
\item \textbf{Neuer offener Punkt \SM:} $\tau_\mu$ direkt aus FFGFT, ohne Umweg über $G_F$ und $v$ --- der einzige konventionsfreie Test des schwachen Sektors.
\end{enumerate}

\section*{Quellen}

Alle Quellen am 5.~September 2026 gegen INSPIRE, APS, Nature, arXiv oder PubMed verifiziert. Zahlenwerte mit exakter rationaler Arithmetik gegen CODATA 2022 berechnet.

\begin{thebibliography}{99}\small
\bibitem{eberhart} M.~Eberhart, M.~Fael, A.~Keshavarzi, D.~Nomura, K.~Schönwald, M.~Steinhauser, T.~Teubner, J.~Wright: \emph{Muon lifetime and Fermi constant: an update}, arXiv:2607.02657 (2026).
\bibitem{mulan} V.~Tishchenko et al.\ (MuLan): Phys.\ Rev.\ D \textbf{87}, 052003 (2013).
\bibitem{sirlin} A.~Sirlin: Phys.\ Rev.\ D \textbf{22}, 971 (1980).
\bibitem{eides26} M.~I.~Eides: \emph{Muonium Hyperfine Splitting Uncertainty Revisited}, arXiv:2510.07281v2 (2026).
\bibitem{liu99} W.~Liu et al.: Phys.\ Rev.\ Lett.\ \textbf{82}, 711 (1999).
\bibitem{eides01} M.~I.~Eides, H.~Grotch, V.~A.~Shelyuto: Phys.\ Rep.\ \textbf{342}, 63 (2001).
\bibitem{morel} L.~Morel, Z.~Yao, P.~Cladé, S.~Guellati-Khélifa: Nature \textbf{588}, 61 (2020).
\bibitem{parker} R.~H.~Parker, C.~Yu, W.~Zhong, B.~Estey, H.~Müller: Science \textbf{360}, 191 (2018).
\bibitem{fan} X.~Fan, T.~G.~Myers, B.~A.~D.~Sukra, G.~Gabrielse: Phys.\ Rev.\ Lett.\ \textbf{130}, 071801 (2023).
\bibitem{aoyama} T.~Aoyama, T.~Kinoshita, M.~Nio: Atoms \textbf{7}, 28 (2019).
\bibitem{besiii} M.~Ablikim et al.\ (BESIII): Phys.\ Rev.\ D \textbf{90}, 012001 (2014).
\bibitem{belle2} I.~Adachi et al.\ (Belle II): Phys.\ Rev.\ D \textbf{108}, 032006 (2023).
\bibitem{kedr} V.~V.~Anashin et al.\ (KEDR): JETP Lett.\ \textbf{85}, 347 (2007).
\bibitem{pdg} PDG 2024, pdgLive S035M: $\tau$ mass average $1776{,}93\pm0{,}09$~MeV.
\bibitem{sturm} S.~Sturm, F.~Köhler, J.~Zatorski et al.: Nature \textbf{506}, 467 (2014).
\bibitem{koehler} F.~Köhler, S.~Sturm, A.~Kracke et al.: J.\ Phys.\ B \textbf{48}, 144032 (2015).
\bibitem{heisse} F.~Heiße et al.: Phys.\ Rev.\ Lett.\ \textbf{131}, 253002 (2023).
\bibitem{li} Q.~Li, C.~Xue, J.-P.~Liu et al.: Nature \textbf{560}, 582 (2018).
\bibitem{rosi} G.~Rosi, F.~Sorrentino, L.~Cacciapuoti, M.~Prevedelli, G.~M.~Tino: Nature \textbf{510}, 518 (2014).
\bibitem{nufit} I.~Esteban, M.~C.~Gonzalez-Garcia, M.~Maltoni, I.~Martinez-Soler, J.~P.~Pinheiro, T.~Schwetz: JHEP \textbf{12} (2024) 216.
\bibitem{codata} P.~J.~Mohr, D.~B.~Newell, B.~N.~Taylor, E.~Tiesinga: Rev.\ Mod.\ Phys.\ \textbf{97}, 025002 (2025).
\end{thebibliography}

\appendix
\section{Betroffene Ableitungsketten}
\label{sec:ketten}

Welche FFGFT-Kette hängt an welchem Vergleichswert, und was ändert der Befund.

\subsection*{Übersicht}

\begin{center}\small
\begin{verbatim}
xi = 4/30000 [K]  (Dok. 338, aus r_mu/r_e und 43200)
      |
  ----+------------------+------------------+
  |                      |                  |
m_e (PDG)         alpha=xi*(E0/MeV)^2   G=xi^2/(4m_e)
[Anker]           (Dok.011,070)         (Dok.A145)
  |                      |                  |
(m_mu/m_e)^2=43200   Vgl: alpha CODATA  Vgl: G Torsionswaage
(Dok.338)            ⚠ Rb/Cs 5sigma     ok modellfrei, sigma=2e-5
  |                  ⚠ HFS ∝ alpha^2
m_mu (pred)
  |
m_tau/m_mu=16.992 (Dok.317)
  |
m_tau (pred)  ─── Vgl: BESIII/BelleII  ok sigma=1e-4
  |
v=m_e/((4/3)*xi^(3/2)) (Dok.006, R104)
  |
G_F=1/(sqrt(2)*v^2)  ─── Vgl: G_F aus tau_mu  ⚠ Delta_q=-0.44%, sqrt(2), 192*pi^3
  |
tau_n, Delta_m_np (Dok.006, R104)  ─── nutzen G_F und alpha
\end{verbatim}
\end{center}

\subsection*{Kette je Befund}

\begin{center}
\begin{tabular}{@{}p{3cm}p{3.5cm}p{3cm}p{4cm}@{}}
\toprule
Befund & Betroffene Dokumente & Zahlenkorrektur & Textkorrektur \\
\midrule
$v$ ist SM-Definition (R112) & Dok.~006, 319, R104 & nein & ,,$v=246$~GeV (gemessen)`` $\to$ ,,(SM-Definition aus $G_F$)`` \\
$G_F$: $\Delta q$, $192\pi^3$, $\sqrt2$ & Dok.~006, R104, 319 & nein & $\Delta q=-0{,}44\,\%$ nennen; mit/ohne $\Delta q$ klären \\
$m_\mu/m_e$ QED-abhängig (R113) & Dok.~338, 317, 011 & nein & ,,modellunabhängig`` $\to$ ,,modellärmst`` \\
$\alpha$ Rb/Cs $5\sigma$ & Dok.~011, 070, 338 & nein & Vergleichswert als CODATA-Mittel mit Spannung zitieren \\
$m_\tau$ zwei Methoden & Dok.~317, 006 & nein & BESIII + Belle II nennen; PDG-Mittel 1776{,}93(9) \\
Neutrinos $\Delta m^2$ & Dok.~339, 340 & nein & NuFIT-Konfiguration einfrieren \\
$G$ & Dok.~A145 & nein & --- \\
\midrule
\textbf{Gesamt} & & \textbf{keine} & Textpräzisierungen \\
\bottomrule
\end{tabular}
\end{center}

\subsection*{Was sich nicht ändert}

Das Residuum $(m_\mu/m_e)^2 = 43200$ vs.\ 42753{,}12 (Rest 446{,}9) ist
$240\,000\,\sigma$ der Messungenauigkeit — physikalisch (R113). Kein Zahlenwert
im Corpus ändert sich. Die Textkorrekturen betreffen ausschließlich die Charakterisierung
der Vergleichswerte sowie den neuen offenen Punkt $\tau_\mu$ direkt aus FFGFT \SM.
